\documentclass[10pt]{article}
\usepackage[margin=20mm]{geometry}
\usepackage{amsmath,amssymb,mathtools,fontspec,unicode-math,xurl,hyperref}
\IfFontExistsTF{DejaVu Serif}{\setmainfont{DejaVu Serif}}{\setmainfont{Latin Modern Roman}}
\setmathfont{latinmodern-math.otf}
\setlength{\parindent}{0pt}\setlength{\parskip}{6pt}\sloppy
\title{Version-limited scientific initial review
M25-0007}
\author{ChatGPT}
\date{25 October 2026}
\begin{document}
\maketitle
Abstract: This report reviews Manjul Bhargava’s Galois groups of random integer polynomials and van der Waerden’s Conjecture, arXiv:2111.06507v3, the exact 34-page supplied manuscript. Six Major findings concern three false statements at their stated scope and three defective derivations of principal results. No counterexample to the named van der Waerden theorem is asserted. Several proof completions and an explicit parameter repair are supplied. The unavailable journal-final text is not adjudicated.\par
=3em\par
\section*{Identity, scope, and principal assessment}
The frozen journal identity is Annals of Mathematics 201 (2025), 339–377, DOI 10.4007/annals.2025.201.2.1. The reviewed version is arXiv v3, submitted 28 September 2024; its PDF front matter displays 1 October 2024. Neither date establishes equivalence to the journal version. The input PDF SHA-256 is\par
00000TOKEN.\par
All 34 pages were read, including references; selected source pages 2–5, 11, 15, 23–24 and 26–30 were rendered and visually inspected. The bound source is one 1464-line `source.tex`; active input and executable-construct checks found none. Page references below are the manuscript’s printed page numbers; source lines are one-based.\par
The central mechanism combines permutation-index bounds, ramification conditions, finite-field Fourier estimates, number-field counting and a three-region argument in conductor/discriminant space. The review checked the definitions, quantifiers, tame/wild distinction, characteristic restrictions, prime-factor splitting, elimination steps, exponents and intransitive convolution. The disposition Major follows the specified standard for a demonstrated invalid derivation of a principal contribution, even where a repair is short.\par
\section*{Problem F01: Theorem 2 includes intransitive groups but counts fields}
Source and location. The definition of \(F_n(G,X)\) on p. 2, Theorem 2 on p. 3 (source lines 165–185), and Section 6, pp. 26–27, especially (23), use a field-counting exponent \(a(G)\) for any permutation group, while \(K_f=\mathbb Q[x]/(f)\) is then an étale algebra.\par
Complete check. Let \(n=3\) and let \(G=\langle(12)\rangle\) fix the third letter. Then \(\operatorname{ind}(G)=1\) and \(u=0\). No degree-three number field has this intransitive permutation Galois group, so the defined \(F_3(G,X)\) is zero. Thus \(a(G)=1/8\) satisfies exactly the stated hypothesis. The second alternative in Theorem 2 would give \[N_3(G,H)=O(H^{(2\cdot3-2)/8+1}\log^2H)
          =O(H^{3/2}\log^2H).\] For \(H\) a positive multiple of four, take \[f_{b,c}(x)=x(x^2+bx+c),\qquad |b|\le H,\quad 2\mid b,
 \quad 0<c\le H,\quad c\equiv2\pmod4.\] The quadratic is Eisenstein at 2, has nonzero discriminant, and has no zero root. The three roots are therefore distinct, and the exact permutation group of \(f_{b,c}\) is the specified \(C_2\) with one fixed point. Its coefficient height is at most \(H\). There are exactly \(H(H+1)/4\) such polynomials. This is \(\Omega(H^2)\), contradicting the asserted bound.\par
Counterevidence and interpretation. Defining \(F_n\) to count étale algebras would remove the vacuous field-counting hypothesis, but it is a different definition and would require corresponding counting inputs. Both p. 2 and Section 6 explicitly say number fields. Theorem 7’s separate treatment of intransitive groups is consistent with the likely intended restriction, but does not restrict Theorem 2 as written.\par
Impact and repair. Major: the advertised general scope of Theorem 2 is false. Restrict it to transitive groups, or formulate and prove an algebra-counting variant. The counterexample does not contradict Theorem 1: its order \(H^2\) is exactly the permitted exceptional order in degree three. Evidence: `independent\_checks.py/json`, exact Eisenstein-family counts.\par
\section*{Problem F02: Corollary 6 omits a necessary degree restriction}
Source and location. Corollary 6 on pp. 4–5, source lines 258–266, asserts \(N_n(G,H)=O(H^{3n/11+1.164})\) for every regular permutation group of order \(n\), citing the exponent \(a(G)=3/8+\epsilon\).\par
Complete check. For \(n=2\), \(G=C_2=S_2\) is genuinely regular. The quadratics \(x^2+bx+c\) just constructed give \(\Omega(H^2)\) distinct polynomials of this group, whereas the claimed exponent is \[\frac6{11}+1.164=\frac{4701}{2750}<2.\] This is a direct counterexample within all the displayed hypotheses.\par
Dependency verification and counterevidence. The published Ellenberg–Venkatesh paper, Proposition 1.3, printed p. 725, explicitly assumes \(n>4\) for this Galois-extension estimate. The older author-hosted 2003 draft encountered in verification has a different lower-degree condition and is not substituted for the published theorem. The input’s reference does not authorize the missing low degrees. High-degree cases are not refuted by the quadratic family.\par
Impact and repair. Major: a stated quantitative corollary is false. Restore the degree restriction required by the actual cited input, and state low-degree bounds separately; degree two needs an exponent at least two. Evidence: the exact counts and the primary published Proposition 1.3 [EV].\par
\section*{Problem F03: Corollary 8 fails for \(S_9\times M_{11}\)}
Source and location. Corollary 8 and its proof, p. 5, source lines 280–284: all \(G\subset S_k\), \(k\le9\), are assigned the bound \(O(H^{8.686})\). The manuscript’s use of \(\subset\) includes equality, explicitly so for \(G=S_n\) after Theorem 2.\par
Construction. Fix a monic integral degree-eleven polynomial \(h\) with Galois group \(M_{11}\) in its natural action. Existence follows from the regular \(M_{11}\) realizations over \(\mathbb Q(t)\) recalled in König’s Section 2 [Koenig] and Hilbert irreducibility. In a resulting \(M_{11}\)-Galois extension of \(\mathbb Q\), the fixed field of a point stabilizer in its natural eleven-point action has degree eleven and Galois closure group \(M_{11}\), since the action is faithful. An integral primitive element of that fixed field gives a monic integral polynomial. This existence theorem is an identified external dependency, not an unverified assertion that one arbitrary specialization works.\par
There is a positive-density coefficient class of monic degree-nine polynomials \(q\) having group \(S_9\). An explicit independently computed class, in ascending coefficient order, is \[(a_0,\ldots,a_9)\equiv
 (323,1594,1589,3841,2809,1924,3976,3358,1162,1)
 \pmod {4199},\] where \(4199=13\cdot17\cdot19\). The three squarefree factorization patterns are respectively \((9)\), \((8,1)\), and \((2,1,1,1,1,1,1,1)\). The exact finite-field tests are saved. Frobenius gives a 9-cycle, hence transitivity, and an 8-cycle fixing a point, hence 2-transitivity. The transposition then forces \(S_9\): the normal subgroup generated by its conjugates is generated by edges of a connected invariant graph, and is the full symmetric group. Thus the coefficient class contains \(\Omega(B^9)\) such \(q\) of height at most \(B\).\par
Let \(L_h,L_q\) be splitting fields. Their intersection is Galois over \(\mathbb Q\). Since \(M_{11}\) is simple, a nontrivial common quotient would be \(M_{11}\) itself; this is impossible for \(S_9\), since \(11\) divides \(|M_{11}|\) and does not divide \(9!\). Thus the fields are linearly disjoint, and \[\operatorname{Gal}(hq)=S_9\times M_{11}.\] The irreducible degrees differ, so the product is squarefree. With the fixed constant \(M_h=\sum|\text{coefficients of }h|\) and \(H(q)\ge1\), coefficient convolution gives \(H(hq)\le M_hH(q)\). Multiplication by \(h\) is injective. Choosing \(B=H/M_h\) proves \[N_{20}(S_9\times M_{11},H)\gg H^9,\] contradicting \(O(H^{8.686})\).\par
Counterevidence, impact, repair. The isolated bound for \(N_{11}(M_{11},H)\) is not contradicted. Theorem 7 itself predicts a maximum of the factor exponents, and the \(S_9\) factor has exponent nine. Major: Corollary 8’s stated range is false. The specific obstruction is removed by taking \(k\le8\), by excluding \(S_9\) when \(k=9\) and supplying the exceptional-group bound, or by allowing exponent nine in this case. No claim is made that merely editing the statement completes every other proof obligation. Evidence: the saved CRT construction and [Koenig].\par
\section*{Problem F04: the stated small parameter does not close the main bound}
Source and location. Section 5, p. 23, chooses \(\delta\) small, suggesting \(1/(2n)\) (source line 895); Case II, p. 24 (line 912 and following display), needs \[E=1+(2+2\delta)(a-u)<n-1\] to absorb its positive logarithmic factor. Here \(a=(n+2)/4\) for \(n\ge6\) and \(u=1/[n(n-1)]\).\par
Check. At \(n=6\), the proposed \(\delta=1/12\) gives \[E=1+\frac{13}{6}\left(2-\frac1{30}\right)
   =\frac{947}{180}>5.\] At \(n=4\), using the stated \(a=1\), it gives \(E=49/16>3\). An \(O(H^E\log^{n-1}H)\) estimate with these exponents does not imply the required \(O(H^{n-1})\).\par
Repair, fully checked. One may instead use \(\delta=1/(4n^2)\) uniformly. For \(n\ge7\), \[E\le\frac{n+4}{2}+\frac{n+2}{8n^2}<n-1,\] and direct rational substitution proves the same strict inequality for \(n=4,5,6\), with \(a=1\) for 4 and 5. For \(n=6\) the sharp requirement here is \(\delta<1/59\). The proposed repair also satisfies the monic Fourier condition \(\delta<1/(2n-1)\), and the squarefull tails retain room for \(\epsilon<\delta/2\). This restores the numerical step in the three-region proof.\par
There is a related dependency mismatch in Section 6, p. 26: \(\delta<1/(n-1)\) is announced and the binary-form Proposition 26 is invoked for monic polynomials. The monic Corollary 33 supplies only \(\delta<1/(2n-1)\). The same smaller common choice fixes this application; the larger stated monic range is not established by that citation.\par
Counterevidence and impact. The freedom to choose \(\delta\) small makes repair straightforward; it does not make the specified numerical value valid. Major under the review standard because this invalid estimate lies directly on the proof of Theorem 1. It is a defective derivation, not a counterexample to Theorem 1. The Case II replacement of \(D^{-u}\) by \(Y^{-u}\) additionally needs partial summation: if \(A(X)\ll X^a\) and \(a>u\), then \(\sum_{D\le Y}D^{-u}\ll Y^{a-u}\), including the integral term. This supplies the estimate; a pointwise inequality would have the wrong direction. The actual exponents used here satisfy \(a>u\).\par
\section*{Problem F05: Theorem 7 differentiates a cumulative upper bound}
Source and location. Section 8, p. 30, source lines 1146–1155, especially line 1153, replaces the number of polynomials of exact height \(H_1\) by \(O(H_1^{\alpha_1-1}\log^{\beta_1}H_1)\) using only the cumulative hypothesis \(N_1(T)=O(T^{\alpha_1}\log^{\beta_1}T)\).\par
Complete logical test. Put \(m(2^j)=2^{2j}\) and \(m(t)=0\) at other positive integers. Then \[\sum_{t\le T}m(t)\le\frac43T^2,
 \qquad \frac{m(2^j)}{2^j}=2^j\longrightarrow\infty.\] Thus a cumulative \(O(T^2)\) bound does not imply the exact-height \(O(T)\) bound used in the proof. This sequence is a counterexample to the inference from the stated hypothesis; it is not claimed to be an actual polynomial counting sequence.\par
Repair in the applications. For the positive power exponents and nonnegative logarithmic exponents used here, a dyadic argument proves the desired convolution estimate without an exact-height estimate. Use \(H_*=\max(1,|\text{coefficients}|)\); fixed-degree comparison with the multiplicative Mahler measure implies \(\prod_i H_*(f_i)\ll H_*(\prod_i f_i)\). Put \(H_*(f_i)\) in dyadic intervals \([2^{j_i},2^{j_i+1})\). Then \(\sum j_i\le L=\log_2(CH)\), up to a harmless fixed adjustment, and the cumulative hypotheses bound the sum by \[\sum_{\substack{j_i\ge0\\\sum j_i\le L}}
 2^{\sum_i\alpha_i j_i}\prod_i(1+j_i)^{\beta_i}.\] If \(\alpha=\max_i\alpha_i>0\), divide by \(2^{\alpha L}\). Every index with \(\alpha_i<\alpha\) has a geometrically decaying weight. The indices with \(\alpha_i=\alpha\) contribute a composition sum of order \(L^{\sum_{\alpha_i=\alpha}(\beta_i+1)-1}\); summing the remaining slack is geometric. This proves the claimed form for these parameter ranges, with upper bounds throughout. Constants in product heights and zero coefficient-height factors must be treated as above. The manuscript does not explicitly restrict arbitrary \(\beta_i\); this repair is not claimed for every negative-logarithm regime.\par
Counterevidence and impact. Theorem 7’s conclusion can survive the invalid proof step in its applications. Major: the step is indispensable to the printed proof of a separately advertised general theorem and to its quantitative applications. The dyadic repair does not rescue F03, since there the maximum exponent is nine. Evidence: exact sparse-shell calculations and the analytic argument just given.\par
\section*{Problem F06: quantitative exponent comparisons fail at stated endpoints}
Source and location. Corollary 3(a)’s proof, pp. 28–29, source lines 1070–1097, uses the power saving \[s(n,k,\epsilon)=\frac{(n-1)(1-1/k)}
 {(n+2)/4-1/[n(n-1)]+\epsilon}.\] Its largest value for \(\epsilon\ge0\) is at zero. The displayed comparisons fail even at that value: \[s(28,5,0)=\frac{81648}{28345}<3,\qquad
 s(53,13,0)=\frac{66144}{18947}<\frac72.\] The second saving is also less than \(7/2\) at \(n=54,55\) and becomes greater at 56. A logarithmic factor only worsens the failed comparisons. The first term \(H^{n+1-k}\) being smaller does not remove the larger second term from their sum.\par
Scope and repair. These are counterexamples to the printed exponent deductions, not number fields refuting the endpoint conclusions. Stronger index information or a separately proved bound could restore the announced thresholds. As sufficient scalar conditions, \(k\ge6,n\ge36\) give \(s> (10/3)(n-1)/(n+2)>3\), and \(k\ge13,n\ge56\) give \(s>(48/13)(n-1)/(n+2)>7/2\). Any use of these repairs must also establish the required group-index assumptions. In the non-elemental branch with \(a=4,k=7\), the denominator furnished by Theorem 2 is \(34/7-u\), not the \(4-u\) printed in (28); its saving is still greater than six at \(n\ge53\), so that branch’s intended conclusion is not refuted by this correction.\par
Related logarithmic comparison. Section 6, p. 28 (source line 1060), asserts \[b=(c-1)-(c-n)\frac{1-1/k}{a+1-1/k-u}<c.\] For the primitive group \(\operatorname{AGL}_2(\mathbb F_3)\) in degree nine, independent enumeration of all 432 affine maps gives index three, so \(c=q(3,6)=3\). Taking the permissible loose field-count exponent \(a=3\) and \(u=1/72\) gives \(b=814/263>3\). The inequality as stated is false. A sharper choice of \(a\) may still imply the advertised theorem bound; this calculation alone does not disprove that bound. Retaining the actual \(b\), or a valid larger logarithmic exponent, avoids this algebraic step.\par
Impact and counterevidence. Major: these invalid comparisons directly support the paper’s advertised quantitative improvements. Other checked thresholds, \(n=10,k=3\) and \(n=16,k=4\), have strictly sufficient savings. No endpoint group counterexample has been inferred from scalar arithmetic. Evidence: exact fractions, the affine-group enumeration, and preserved failed check attempts.\par
\section*{Positive checks, resolved concerns, and dependency limits}
\subsection*{Group theory and ramification}
The transposition argument closes after adding reflexivity to the equivalence relation; conjugation supplies transitivity. For a nontrivial inertia element, the orbit count gives the wild discriminant lower bound; tame primes \(p>n\) give equality for the inertia generator. Small primes form a fixed finite set with bounded degree-dependent discriminant exponents.\par
Lemma 24, p. 15, prints \(v_p(D)=ef\) and assigns index \(e\) to \(x^e\). The correct quantities are \((e-1)f\) and \(e-1\). Over an unramified splitting extension, the \(f\) conjugate factors have reduction \((x-c_i)^e\), so their total index is at least \(f(e-1)\). Translation signs can likewise be corrected. These paired local errors do not refute the lemma’s needed statement and are not separately inflated into an additional Major finding.\par
The induction for the movement inequality (6), p. 11, omits the fully paired odd-\(k\) case and boundary cases where its induction hypothesis changes range. An independent completion avoids those cases. Write \(Y=\lfloor m/2\rfloor\) and start with \(Y\) disjoint transpositions. Choose \(y\) of them uniformly. Every \(k\)-set moved by the full family is hit by at least one asymmetric pair; the probability that a chosen pair hits is at least \(y/Y\). Thus \(f(m,k,y)\ge(y/Y)f(m,k,Y)\). For even \(m\), odd \(k\) has no fixed \(k\)-sets; for even \(k\ge2\), the fixed proportion \(\binom{Y}{k/2}/\binom{2Y}{k}\) is at most \(1/(2Y-1)\le1/5\) in the relevant range, by the ratio between successive even \(k\). For odd \(m\), the fixed proportion is at most \(1/m\), by the analogous parity-wise ratios starting at \(k=1,2\). These bounds imply the required \[f(m,k,y)\ge \frac{8y}{5k}\binom{m-1}{k-1}.\] The small allowed cases are immediate. Exact enumeration checked 17,545 parameter triples with \(3\le m\le60\); the universal argument, rather than this finite test, supplies the completion.\par
\subsection*{Geometry, Fourier estimates, and the three regions}
The Vandermonde argument in Proposition 22 uses the no-zero-subset-sum assumption to exclude solutions at infinity and gives the stated finite intersection bound. In Proposition 21, \(p>n\) ensures that positive multiplicities and their subset sums remain nonzero and that Newton identities can be inverted. The monic-to-binary normalization in Remark 23 needs the scalar factor \(t^{-n}\); with it the translation/dilation orbit argument works. The zero binary form must be removed before asserting a bounded factorization multiplicity; nonzero points retain the requisite density.\par
For the monic Fourier bound, the highest nonzero dual coordinate supplies a nonconstant translation-phase polynomial of degree at most \(n\); its leading binomial coefficient is nonzero for \(p>n\). The square-root bound and the CRT strata then have the necessary hypotheses. The distinction between binary and monic exponents is retained in F04.\par
In Case III the double discriminant is a nonzero polynomial: the family \(x^n+bx+c\) supplies a direct nonvanishing specialization. Its zero locus gives \(O(H^{n-2})\) choices before the last coefficient. Large prime factors have bounded number because the double discriminant has polynomial height and each such prime exceeds a fixed positive power of \(H\). In the other subcase, a greedy product of primes at most \(H^{\delta/2}\) crosses \(H^{1+\delta/2}\) without exceeding \(H^{1+\delta}\). The powerful-integer tails have strict power savings after the repaired choice of \(\delta\).\par
Section 6’s triangular elimination needs more than dimension alone. The earlier weighted moment equations give a finite projection to the first \(n-k\) coordinates; integrality supplies monic annihilating relations for the remaining coordinates. This justifies the bounded intermediate fibers required by the counting argument. Theorem 20’s trace-zero lattice count and the primitive-extension condition were checked together: every nonrational element generates in that setting, and the discriminant divisibility constraints fit the successive coefficient ranges.\par
\subsection*{What this review does not establish}
The manuscript uses classification-based group-index estimates, classical number-field counting, geometry-of-numbers results, Weil bounds and Hilbert irreducibility as external theorems. Their stated interfaces and relevant restrictions were checked; this report is not a fresh proof of the whole cited literature. No author executable or computational output was run or adopted as independent verification. The special hyperelliptic application was read with its genus/degree conventions; no unqualified low-genus objection is asserted. Finite computations test the explicit algebra above and do not certify all degrees or all proofs.\par
\section*{Conclusion and version disposition}
ID | Judgment | Repair/adoption | Remaining work\par
F01 | Major; false scope | Transitive restriction | Algebra variant needs new hypotheses\par
F02 | Major; false at \(n=2\) | Restore degree range | Separate low-degree corollaries\par
F03 | Major; \(k=9\) false | Exclude \(S_9\) or exponent 9 | Re-state valid product range\par
F04 | Major; invalid estimate | Smaller \(\delta\) checked | Edit dependent monic citations\par
F05 | Major; invalid inference | Dyadic proof for applications | State parameter range precisely\par
F06 | Major; invalid comparisons | Exact failures retained | Strengthen endpoint input or weaken thresholds\par
None of these judgments is transferred to a journal text not supplied. The arXiv record inspected lists v3 as the latest version. The publisher records receipt on 14 January 2022, revision on 13 February 2024, acceptance on 14 February 2024 and online publication on 12 March 2025. These dates do not prove textual identity. Bounded primary-source searches did not identify an article-specific erratum; that negative search is not a universal absence claim. The separately titled arXiv:2410.03792 is an exposition, not a v4 or an erratum, and was not treated as a replacement.\par
\section*{Supporting checks}
The task evidence directory is `verification/0007`. It contains the preserved source-text extraction, static source audit, primary-source web records, exact check source/results, the version record, coverage, findings, provenance, build receipts and report-page inspection records. The primary check script SHA-256 is\par
00001TOKEN.\par
The executed checks include six Eisenstein-family sizes, 97 rational parameter rows, 18 sparse-shell rows, 17,545 movement cases, exact finite-field irreducibility/CRT checks, endpoint fractions and the 432-element affine enumeration. The first two check attempts genuinely failed because the draft expectations incorrectly anticipated passing endpoint inequalities; their sources and failure records are preserved. The corrected script records the failures of the manuscript’s comparisons as its expected mathematical result and exited zero. This distinction prevents a passing script from being represented as a passing paper.\par
The reproducible scientific commands are Python invocations of `independent\_checks.py` and `additional\_checks.py`. The build uses the previously authorized Tectonic 0.17.0 with `–untrusted –keep-logs –keep-intermediates –synctex –makefile-rules`; shell escape and author-code execution are disabled. Actual stdout/stderr, return code, source before/after hashes, compiler hash, dependency hashes and PDF hash are retained outside the scientific narrative. The exact final hashes and every inspected page are bound by the delivery records; this paragraph does not substitute for those records.\par
The controller’s observed account and webpage model are attributed to the controller. The assistant self-identification is stored separately. No backend-equivalence claim is made. All writes and bounded runtime caches belong to this task staging.\par
9 M. Bhargava, Galois groups of random integer polynomials and van der Waerden’s Conjecture, arXiv:2111.06507v3. https://arxiv.org/abs/2111.06507v3. Journal identity and editorial dates. https://annals.math.princeton.edu/2025/201-2/p01. J. S. Ellenberg and A. Venkatesh, The number of extensions of a number field with fixed degree and bounded discriminant, Ann. of Math. 163 (2006), 723–741, Proposition 1.3. https://annals.math.princeton.edu/wp-content/uploads/annals-v163-n2-p11.pdf. J. König, On rational functions with monodromy group \(M_{11}\), arXiv:1503.06686v2, Section 2. https://arxiv.org/pdf/1503.06686v2.\par
\end{document}
